% !TeX root = ../教材题库_全部条目_转写工作稿.tex
% Source: XJ-B1-C06-TH-P218-Y0289
% 第6章 / 6.2 抽样 / 教材 p.218 / PDF 第225页
% OCR未校对：公式与文字均可能有误；未完成数学TeX转写；禁止直接发给学生
% 原题及相邻原解截图（核对用）：../题目原图/XJ-B1-C06-TH-P218-Y0289-01.webp
% 本文件只在转写完成后作为题干使用；OCR 条目尚不是可靠题干。
\begin{bankitem}{p.218 思考（含紧邻说明）}{第6章 / 6.2 抽样 / 教材 p.218 / PDF 第225页}{XJ-B1-C06-TH-P218-Y0289}{OCR未校对：公式与文字均可能有误；未完成数学TeX转写；禁止直接发给学生}
\begin{ocrdraft}
[教材p.218 / PDF第225页]
\par
？思考
\par
如果要通过抽样了解一所学校高一年级学生的身高情况，你认为应该
\par
怎样抽样？
\par
我们知道，影响身高的因素有很多，性别是其中一个主要因素．高中男生的身
\par
高普遍高于女生的身高，若用简单随机抽样的方法抽取样本，则得到的结果可能会
\par
比较“极端”．因而需要学习一种新的抽样方法．
\par
当总体由差异明显的几个部分组成时，为了使抽取的样本更好地反映总体的情
\par
况，把总体中各个个体按照某种特征或某种规则划分为互不交叉的层，然后对各层
\par
进行简单随机抽样，这种抽样方法称为分层抽样．
\par
对于上述“思考”中的问题，由于高中男生与女生的身高情况差异较为明显，
\par
而相同性别的身高差异相对较小．于是，我们可以利用性别和身高的这种关系，把
\par
高一年级学生按男生和女生分为两层，再对各层分别进行简单随机抽样，那么每层
\par
的样本量应如何分配？
\par
自然地，为了使样本的结构与总体的分布相近，人数多的群体应多抽一些，人
\par
数少的群体应少抽一些．因此，按男生、女生在全体高一年级学生中所占的比例进
\par
行分配是一种比较合理的方法，即
\par
男生人数
\par
女生人数
\par
男生样本量一
\par
\ensuremath{\times}总样本量，女生样本量一
\par
\ensuremath{\times}总样本量
\par
高一学生数
\par
高一学生数
\par
这样无论是男生还是女生，每个学生被抽到的概率都相等，并且有效地避免了
\par
“极端”样本．
\par
在实际抽样调查中，样本量的分配是一个比较复杂的问题，要充分考虑诸多要
\par
素，如估计精度、调查费用等，再决定分配方案．
\end{ocrdraft}
\end{bankitem}
